\documentclass[11pt,a4paper,openany]{report}

% Fichier autonome : téléverser uniquement synthese.tex dans Overleaf.
\usepackage[T1]{fontenc}
\usepackage[utf8]{inputenc}
\usepackage[french]{babel}
\usepackage{lmodern}
\usepackage{microtype}
\usepackage[a4paper,margin=24mm,headheight=15pt]{geometry}
\usepackage{xcolor}
\usepackage{booktabs}
\usepackage{longtable}
\usepackage{enumitem}
\usepackage{amsmath}
\usepackage{listings}
\usepackage[most]{tcolorbox}
\usepackage{fancyhdr}
\usepackage{hyperref}
\usepackage{bookmark}

\definecolor{GuideBlue}{HTML}{174A7E}
\definecolor{GuideGreen}{HTML}{237A57}
\definecolor{GuideOrange}{HTML}{B45F06}
\definecolor{GuideRed}{HTML}{A61B1B}
\definecolor{GuideGray}{HTML}{F4F5F7}

\hypersetup{
  colorlinks=true,
  linkcolor=GuideBlue,
  urlcolor=GuideBlue,
  pdftitle={Synthèse logique avec Genus},
  pdfauthor={Karim Sabra}
}

\pagestyle{fancy}
\fancyhf{}
\fancyhead[R]{\small\nouppercase{\leftmark}}
\fancyfoot[C]{\thepage}
\renewcommand{\headrulewidth}{0.4pt}
\setcounter{tocdepth}{2}
\setlist{nosep,leftmargin=*}

\lstdefinelanguage{SystemVerilog}{
  morekeywords={assign,begin,end,if,else,for,logic,module,endmodule,input,
    output,initial,task,endtask,automatic,localparam,parameter,int,unsigned,
    always,posedge,negedge},
  sensitive=true,
  morecomment=[l]{//},
  morecomment=[s]{/*}{*/},
  morestring=[b]"
}

\lstdefinelanguage{TclGuide}{
  morekeywords={set,proc,if,elseif,else,foreach,return,error,catch,puts,exit,
    source,global,file,info,dict,string,read_hdl,elaborate,read_sdc,
    check_design,check_timing_intent,syn_generic,syn_map,syn_opt,
    report_timing,report_area,report_qor,write_hdl,write_sdc,read_mmmc,
    read_physical,init_design,create_library_set,create_rc_corner,
    create_delay_corner,create_constraint_mode,create_analysis_view,
    set_analysis_view},
  sensitive=true,
  morecomment=[l]{\#},
  morestring=[b]"
}

\lstdefinestyle{guidecode}{
  basicstyle=\ttfamily\scriptsize,
  keywordstyle=\color{GuideBlue}\bfseries,
  commentstyle=\color{GuideGreen},
  stringstyle=\color{GuideOrange},
  backgroundcolor=\color{GuideGray},
  frame=single,
  rulecolor=\color{black!20},
  columns=fullflexible,
  keepspaces=true,
  showstringspaces=false,
  breaklines=true,
  tabsize=2,
  numbers=left,
  numberstyle=\tiny\color{black!50},
  xleftmargin=1.5em,
  literate=
    {à}{{\`a}}1 {ç}{{\c c}}1 {é}{{\'e}}1 {è}{{\`e}}1
    {ê}{{\^e}}1 {î}{{\^i}}1 {ô}{{\^o}}1 {ù}{{\`u}}1
}

\newtcolorbox{keybox}[1][À retenir]{
  breakable,colback=GuideBlue!5,colframe=GuideBlue,title=\textbf{#1}}
\newtcolorbox{warningbox}[1][Attention]{
  breakable,colback=GuideOrange!7,colframe=GuideOrange,title=\textbf{#1}}
\newtcolorbox{failurebox}[1][Diagnostic]{
  breakable,colback=GuideRed!5,colframe=GuideRed,title=\textbf{#1}}

\begin{document}

\begin{titlepage}
  \centering
  \vspace*{28mm}
  {\Huge\bfseries\color{GuideBlue}Synthèse logique avec Genus\par}
  \vspace{7mm}
  {\LARGE Du RTL à une netlist mappée réutilisable\par}
  \vfill
  {\Large Karim Sabra\par}
  \vspace{4mm}
  {Guide autonome compatible Overleaf\par}
  {\today\par}
\end{titlepage}

\chapter*{Mode d'emploi Overleaf}
\addcontentsline{toc}{chapter}{Mode d'emploi Overleaf}

Ce document ne dépend d'aucun autre fichier LaTeX. Dans Overleaf, créer un
projet vide, téléverser uniquement \texttt{synthese.tex}, le sélectionner
comme document principal et choisir le compilateur pdfLaTeX.

\begin{warningbox}[Portée]
Les noms des bibliothèques, des corners et des fichiers technologiques sont
des exemples. Il faut employer les chemins, unités et cellules autorisés par
le laboratoire. Un processus Genus terminé sans erreur ne prouve pas à lui
seul que le timing ou les contraintes sont corrects.
\end{warningbox}

\tableofcontents

\chapter{But et progression}

\section{Ce que produit la synthèse}

Genus transforme un RTL élaboré en une netlist composée de cellules de la
bibliothèque cible. Il optimise la logique sous les contraintes de timing,
d'aire, de transition, de fanout et de charge.

\begin{longtable}{p{0.20\textwidth}p{0.70\textwidth}}
\toprule
Niveau & Objectif \\
\midrule
1 & Lire le full adder, l'élaborer et comprendre le SDC. \\
2 & Exécuter les trois phases \texttt{syn\_generic}, \texttt{syn\_map} et
\texttt{syn\_opt}. \\
3 & Générer les rapports et les deux fichiers nécessaires à Innovus. \\
4 & Ajouter un wrapper, le contrôle des entrées, les vues MMMC et les données
physiques. \\
\bottomrule
\end{longtable}

\section{Entrées et sorties}

\begin{longtable}{p{0.23\textwidth}p{0.30\textwidth}p{0.37\textwidth}}
\toprule
Élément & Exemple & Rôle \\
\midrule
RTL & \texttt{full\_adder\_comb.sv} & description logique synthétisable \\
Filelist & \texttt{rtl.f} & ordre et chemins des sources \\
Contraintes & \texttt{constraints.sdc} & clocks, délais d'E/S et exceptions \\
Liberty & \texttt{stdcells\_tc.lib} & fonctions, arcs, délais et puissance \\
Netlist mappée & \texttt{full\_adder\_comb.mapped.v} & cellules choisies \\
SDC exporté & \texttt{mapped.sdc} & contraintes transmises \\
Rapports & timing, aire, QoR & preuves à examiner \\
\bottomrule
\end{longtable}

\chapter{Exemple minimal : full adder}

\section{RTL}

\begin{lstlisting}[style=guidecode,language=SystemVerilog,
caption={rtl/full\_adder\_comb.sv}]
module full_adder_comb (
    input  logic a_i,
    input  logic b_i,
    input  logic cin_i,
    output logic sum_o,
    output logic cout_o
);
    assign sum_o  = a_i ^ b_i ^ cin_i;
    assign cout_o = (a_i & b_i)
                  | (a_i & cin_i)
                  | (b_i & cin_i);
endmodule
\end{lstlisting}

\section{Filelist}

\begin{lstlisting}[style=guidecode,caption={syn/rtl.f}]
rtl/full_adder_comb.sv
\end{lstlisting}

Les chemins relatifs doivent être interprétés depuis un dossier connu. Le
script Genus ci-dessous se place explicitement à la racine de travail avant de
lire la filelist.

\section{SDC combinatoire}

Un circuit sans clock peut être contraint par une clock virtuelle. Elle sert
de référence aux délais d'entrée et de sortie.

\begin{lstlisting}[style=guidecode,language=TclGuide,
caption={syn/constraints.sdc}]
set PERIOD_NS 2.000
create_clock -name VCLK -period $PERIOD_NS

set_input_delay  0.200 -clock VCLK [get_ports {a_i b_i cin_i}]
set_output_delay 0.200 -clock VCLK [get_ports {sum_o cout_o}]

set_input_transition 0.050 [get_ports {a_i b_i cin_i}]
set_load 0.010 [get_ports {sum_o cout_o}]
\end{lstlisting}

\begin{keybox}[Interprétation]
Le budget disponible pour un chemin entrée-vers-sortie est approximativement
la période moins le délai d'entrée, le délai de sortie et les marges internes.
Les valeurs doivent correspondre au contrat réel du bloc, pas à un nombre
choisi uniquement pour obtenir un rapport positif.
\end{keybox}

\chapter{Commande et Tcl minimal}

\section{Commande directe}

Avec un fichier Tcl nommé \texttt{genus\_minimal.tcl} :

\begin{lstlisting}[style=guidecode,language=bash]
export LIBERTY_TC=/chemin/autorise/stdcells_tc.lib
genus -files genus_minimal.tcl
\end{lstlisting}

La variable d'environnement évite d'inscrire un chemin privé dans le script
partagé.

\section{Script complet minimal}

\begin{lstlisting}[style=guidecode,language=TclGuide,
caption={genus\_minimal.tcl}]
proc require_env {name} {
    if {![info exists ::env($name)] ||
        [string trim $::env($name)] eq ""} {
        error "variable obligatoire absente: $name"
    }
    return $::env($name)
}

set TOP full_adder_comb
set ROOT [file normalize .]
set FILELIST [file join $ROOT syn rtl.f]
set SDC [file join $ROOT syn constraints.sdc]
set REPORTS [file join $ROOT reports]
set OUTPUTS [file join $ROOT outputs]
file mkdir $REPORTS
file mkdir $OUTPUTS

set LIB_TC [file normalize [require_env LIBERTY_TC]]
if {![file exists $LIB_TC] || [file size $LIB_TC] == 0} {
    error "Liberty absent ou vide: $LIB_TC"
}

set_db hdl_language sv
set_db library [list $LIB_TC]
set_db init_hdl_search_path [list $ROOT]
cd $ROOT

read_hdl -sv -f $FILELIST
elaborate $TOP
read_sdc $SDC

redirect -file [file join $REPORTS check_design.rpt] {
    check_design -all
}
redirect -file [file join $REPORTS check_timing_intent.rpt] {
    check_timing_intent -verbose
}

syn_generic
syn_map
syn_opt

redirect -file [file join $REPORTS report_timing.rpt] {
    report_timing -max_paths 20
}
redirect -file [file join $REPORTS report_area.rpt] {
    report_area
}
redirect -file [file join $REPORTS report_qor.rpt] {
    report_qor
}

set NETLIST [file join $OUTPUTS ${TOP}.mapped.v]
set SDC_OUT [file join $OUTPUTS ${TOP}.mapped.sdc]
write_hdl > $NETLIST
write_sdc > $SDC_OUT

foreach path [list $NETLIST $SDC_OUT] {
    if {![file exists $path] || [file size $path] == 0} {
        error "sortie absente ou vide: $path"
    }
}

puts "GENUS_TUTORIAL_STATUS: BASIC_FLOW_COMPLETED"
exit 0
\end{lstlisting}

\section{Lecture étape par étape}

\begin{enumerate}
  \item \texttt{read\_hdl} compile les sources.
  \item \texttt{elaborate} construit la hiérarchie et résout les paramètres.
  \item \texttt{read\_sdc} applique l'intention de timing.
  \item Les checks détectent les références non résolues et contraintes
    incomplètes avant l'optimisation.
  \item \texttt{syn\_generic} optimise une représentation indépendante de la
    technologie.
  \item \texttt{syn\_map} choisit les cellules Liberty.
  \item \texttt{syn\_opt} améliore le design mappé sous contraintes.
  \item Les rapports qualifient le résultat; les exports alimentent Innovus.
\end{enumerate}

\chapter{Comprendre les rapports}

\section{Timing}

Pour un check de setup :
\[
  \mathrm{slack} = t_{\mathrm{requis}} - t_{\mathrm{arrivee}}.
\]
Un slack négatif indique une violation. Le WNS est le pire slack et le TNS la
somme des slacks négatifs. Il faut aussi vérifier :

\begin{itemize}
  \item que les clocks et les unités sont celles attendues;
  \item qu'aucun chemin important n'est non contraint;
  \item que les ports reçoivent bien délais, transitions et charges;
  \item que les exceptions ciblent les objets voulus;
  \item que les violations de transition, capacitance et fanout sont séparées
    des violations de timing.
\end{itemize}

\section{Aire et QoR}

L'aire totale doit être lue avec le nombre de cellules, la répartition
combinatoire/séquentielle et les éventuelles cellules non mappées. Une aire
plus petite n'est pas automatiquement meilleure si elle dégrade le timing ou
la robustesse électrique.

\begin{warningbox}[Verdict correct]
Le code retour de l'outil, l'élaboration, les contraintes, le timing, les
design rules et les exports sont des contrôles distincts. Ne pas les résumer
par un seul message « terminé ».
\end{warningbox}

\chapter{Wrapper réutilisable}

Le wrapper suivant contrôle les fichiers avant de lancer Genus, isole les
résultats et conserve le code retour.

\begin{lstlisting}[style=guidecode,language=bash,
caption={run\_genus.sh}]
#!/usr/bin/env bash
set -Eeuo pipefail

GENUS_BIN=${GENUS_BIN:-genus}
RUN_ROOT=${RUN_ROOT:-/tmp/digi_tuto_runs}
RUN_ID=${RUN_ID:-$(date -u +%Y%m%dT%H%M%SZ)}
RUN_DIR="$RUN_ROOT/genus/full_adder_comb/$RUN_ID"

: "${LIBERTY_TC:?définir LIBERTY_TC}"
for path in \
  "$LIBERTY_TC" \
  rtl/full_adder_comb.sv \
  syn/rtl.f \
  syn/constraints.sdc \
  genus_minimal.tcl; do
  [[ -s "$path" ]] || {
    echo "ERREUR: fichier absent ou vide: $path" >&2
    exit 2
  }
done

command -v "$GENUS_BIN" >/dev/null 2>&1 || {
  echo "ERREUR: Genus introuvable" >&2
  exit 127
}

[[ ! -e "$RUN_DIR" ]] || {
  echo "ERREUR: dossier déjà présent: $RUN_DIR" >&2
  exit 2
}
mkdir -p "$RUN_DIR"

set +e
"$GENUS_BIN" -files genus_minimal.tcl \
  -log "$RUN_DIR/genus.log" \
  2>&1 | tee "$RUN_DIR/console.log"
RC_GENUS=${PIPESTATUS[0]}
RC_TEE=${PIPESTATUS[1]}
set -e

(( RC_GENUS == 0 && RC_TEE == 0 )) || {
  echo "RESULTAT: FAIL (genus=$RC_GENUS tee=$RC_TEE)" >&2
  exit 1
}

for path in \
  outputs/full_adder_comb.mapped.v \
  outputs/full_adder_comb.mapped.sdc \
  reports/report_timing.rpt; do
  [[ -s "$path" ]] || {
    echo "RESULTAT: FAIL (sortie absente: $path)" >&2
    exit 3
  }
done

echo "RESULTAT: FLOW_EXECUTE"
echo "Lire reports/report_timing.rpt avant le handoff"
\end{lstlisting}

\begin{keybox}
Dans un flow de laboratoire, il est préférable de placer aussi les dossiers
\texttt{reports} et \texttt{outputs} sous \texttt{RUN\_DIR}. Le script minimal
les garde près du Tcl pour rester lisible; le wrapper complet fourni avec le
tutoriel les isole sous \texttt{/tmp/digi\_tuto\_runs}.
\end{keybox}

\chapter{Méthodologie avancée}

\section{Passer du TC au MMMC}

Une analyse MMMC sépare :

\begin{itemize}
  \item les \emph{library sets} Liberty;
  \item les corners RC;
  \item les delay corners qui associent timing et interconnexions;
  \item les modes de contraintes;
  \item les analysis views pour setup et hold.
\end{itemize}

\begin{lstlisting}[style=guidecode,language=TclGuide,
caption={Structure MMMC simplifiée}]
create_library_set -name LIB_BC -timing [list $LIBERTY_BC]
create_library_set -name LIB_TC -timing [list $LIBERTY_TC]
create_library_set -name LIB_WC -timing [list $LIBERTY_WC]

create_rc_corner -name RC_BC -qx_tech_file $QRC_BC
create_rc_corner -name RC_TC -qx_tech_file $QRC_TC
create_rc_corner -name RC_WC -qx_tech_file $QRC_WC

create_delay_corner -name DC_BC \
    -library_set LIB_BC -rc_corner RC_BC
create_delay_corner -name DC_TC \
    -library_set LIB_TC -rc_corner RC_TC
create_delay_corner -name DC_WC \
    -library_set LIB_WC -rc_corner RC_WC

create_constraint_mode -name FUNC -sdc_files [list $SDC]

create_analysis_view -name VIEW_BC \
    -constraint_mode FUNC -delay_corner DC_BC
create_analysis_view -name VIEW_TC \
    -constraint_mode FUNC -delay_corner DC_TC
create_analysis_view -name VIEW_WC \
    -constraint_mode FUNC -delay_corner DC_WC

set_analysis_view \
    -setup [list VIEW_TC VIEW_WC] \
    -hold  [list VIEW_TC VIEW_BC]
\end{lstlisting}

Le choix exact des vues setup et hold appartient à la méthodologie du PDK. Il
ne faut pas déduire ce choix uniquement du nom « best » ou « worst ».

\section{Données physiques}

Une synthèse physical-aware ajoute les LEF technologiques et de cellules ainsi
que les fichiers QRC :

\begin{lstlisting}[style=guidecode,language=TclGuide]
read_mmmc flow_mmmc.tcl
read_physical -lef [list $TECH_LEF $STDCELL_LEF]
read_hdl -sv -f syn/rtl.f
elaborate $TOP
init_design
\end{lstlisting}

Cette estimation améliore la corrélation avec Innovus, mais ne remplace pas
placement, CTS, routage et extraction réelle.

\section{Design séquentiel}

Pour un additionneur pipeliné, le SDC doit définir la vraie clock :

\begin{lstlisting}[style=guidecode,language=TclGuide]
create_clock -name CLK -period 2.000 [get_ports clk_i]
set_clock_uncertainty 0.100 [get_clocks CLK]
set_input_delay  0.200 -clock CLK \
    [remove_from_collection [all_inputs] [get_ports clk_i]]
set_output_delay 0.200 -clock CLK [all_outputs]
set_false_path -from [get_ports rst_ni]
\end{lstlisting}

L'exception de reset doit correspondre à l'architecture et à la politique du
laboratoire. Vérifier les objets ciblés avec les rapports de contraintes.

\chapter{Handoff vers Innovus}

Le package minimal contient :

\begin{enumerate}
  \item la netlist mappée non vide;
  \item le SDC exporté non vide;
  \item le nom exact du top;
  \item les mêmes familles Liberty utilisées par la configuration Innovus;
  \item les rapports de checks, de timing et de QoR;
  \item une liste explicite des violations ou hypothèses restantes.
\end{enumerate}

\begin{failurebox}[Gate avant PnR]
Ne pas lancer Innovus si la netlist contient des références non résolues, si
le SDC n'a pas produit les clocks et contraintes attendues, ou si le choix des
corners est inconnu. Un timing positif au seul corner typique reste un résultat
typique, pas une preuve MMMC.
\end{failurebox}

\chapter{Diagnostic}

\begin{longtable}{p{0.23\textwidth}p{0.31\textwidth}p{0.36\textwidth}}
\toprule
Symptôme & Cause probable & Première action \\
\midrule
Genus introuvable & environnement Cadence non chargé &
Vérifier \texttt{command -v genus}. \\
Module non résolu & filelist ou ordre incorrect &
Lire la première erreur de \texttt{read\_hdl}. \\
Top inconnu & mauvais nom d'élaboration &
Comparer le module RTL et la variable \texttt{TOP}. \\
Clock absente & SDC non lu ou objet vide &
Examiner \texttt{check\_timing\_intent} et les clocks. \\
Chemins non contraints & contrat d'E/S incomplet &
Ajouter les contraintes justifiées, sans masquer les chemins. \\
Cellules non mappées & Liberty incompatible ou incomplète &
Vérifier la library active et le rapport de mapping. \\
WNS négatif & chemin setup trop lent &
Lire le premier chemin complet avant de changer l'effort. \\
Export vide & étape interrompue ou mauvais chemin &
Contrôler le code retour et la taille de chaque fichier. \\
\bottomrule
\end{longtable}

\appendix
\chapter{Checklist finale}

\begin{enumerate}
  \item Toutes les sources de la filelist existent et le top est explicite.
  \item L'élaboration ne contient pas de référence non résolue.
  \item Les unités Liberty et SDC sont cohérentes.
  \item Les clocks, délais d'E/S, transitions et charges sont vérifiés.
  \item Les chemins non contraints et les exceptions sont examinés.
  \item Les rapports timing, aire, QoR et design rules sont non vides.
  \item La netlist mappée et le SDC exporté sont non vides.
  \item Les résultats TC et MMMC restent clairement distingués.
  \item Le package Innovus reprend le top et les corners corrects.
\end{enumerate}

\end{document}
